\subsection{Minlos's lemma}

Let T be a finite-dimensional vector space and \(\mu\) a bounded measure on \(T'\) ; we shall identify T with the dual of \(T'\) , so that the Fourier transform \(\Phi\) of \(\mu\) is a function on T. We assume given two positive quadratic forms h and q on T and a number \(\varepsilon > 0\) . For every real number r > 0, we denote by \(C_{r}\) the set of \(x' \in T'\) such that \(\langle x, x' \rangle^{2} \leqslant r^{2} h(x)\) for all \(x \in T\) .

PROPOSITION 10. --- Under the hypothesis \(\Phi(0) - \mathcal{R}\Phi \leqslant \varepsilon + q\) , we have

\[
\mu (\mathrm{T} ^ {\prime} - \mathrm{C} _ {r}) \leqslant 3 \left(\varepsilon + r ^ {- 2} \operatorname{Tr} (q / h)\right)\tag{55}
\]

for every \(r > 0\) .

One writes \(\operatorname{Tr}(q/h)\) for the trace of q with respect to h (cf.~Annex, No.~1). The formula (55) is trivial when \(\operatorname{Tr}(q/h)\) is infinite. We assume henceforth that \(\operatorname{Tr}(q/h)\) is finite, hence that \(h(x)=0\) implies \(q(x)=0\) for \(x\in T\) .

Let \(a_1, \ldots, a_n\) be elements of \(\mathbf{T}\) , and \(\mathbf{D}\) the set of \(x' \in \mathbf{T}'\) such that \(\sum_{j=1}^{n} \langle a_j, x' \rangle^2 > 1\) . For every real \(t \geqslant 0\) we have \(3(1 - e^{-t/2}) \geqslant 0\) , and we even have

\[
3 (1 - e ^ {- t / 2}) \geqslant 3 (1 - e ^ {- 1 / 2}) \geqslant 3 \left(1 - \left(\frac {9}{4}\right) ^ {- 1 / 2}\right) = 1
\]

for \(t > 1\) , because \(e > \frac{9}{4}\) . Applying these inequalities to \(t = \sum_{j=1}^{n} \langle a_j, x' \rangle^2\) , we obtain

\[
\mu (\mathrm{D}) \leqslant 3 \int_ {\mathrm{T} ^ {\prime}} \left(1 - \exp \left(- \frac {1}{2} \sum_ {j = 1} ^ {n} \langle a _ {j}, x ^ {\prime} \rangle^ {2}\right)\right) d \mu (x ^ {\prime}).\tag{56}
\]

Let \(\gamma\) be the measure on \(\mathbf{R}\) having density \(t \mapsto (2\pi)^{-1/2} e^{-t^2/2}\) with respect to Lebesgue measure. By Lemma 2 of No.~4,

\[
\int_ {\mathbf {R}} e ^ {i u t} d \gamma (t) = e ^ {- u ^ {2} / 2}
\]

for all real \(u\) . Consequently,

\[
\begin{array}{l} 1 - \exp \left(- \frac {1}{2} \sum_ {j = 1} ^ {n} \langle a _ {j}, x ^ {\prime} \rangle^ {2}\right) \\ = \int \dots \int \left(1 - e ^ {i \sum_ {j = 1} ^ {n} \langle a _ {j}, x ^ {\prime} \rangle t _ {j}}\right) d \gamma (t _ {1}) \dots d \gamma (t _ {n}) \end{array}\tag{57}
\]

for all \(x' \in T'\) . The function of \(x', t_{1}, \ldots, t_{n}\) to be integrated in the second member is continuous and is bounded above in absolute value by 2, and the measures \(\mu\) and \(\gamma\) are bounded; one can therefore integrate the two members of (57) with respect to \(d\mu(x')\) and interchange the integrations with respect to \(\mu\) and \(\gamma\) ; one obtains

\[
\begin{array}{l} \int_ {\mathbf {T} ^ {\prime}} \left(1 - \exp \left(- \frac {1}{2} \sum_ {j = 1} ^ {n} \langle a _ {j}, x ^ {\prime} \rangle^ {2}\right)\right) d \mu (x ^ {\prime}) \\ = \int \dots \int \left(\Phi (0) - \Phi \left(\sum_ {j = 1} ^ {n} t _ {j} a _ {j}\right)\right) d \gamma (t _ {1}) \dots d \gamma (t _ {n}). \end{array}\tag{58}
\]

Since \(q\) is a quadratic form on \(\mathbf{T}\) , there exist real numbers \(q_{jk}\) such that

\[
q \left(\sum_ {j = 1} ^ {n} t _ {j} a _ {j}\right) = \sum_ {j, k} q _ {j k} t _ {j} t _ {k}
\]

for \(t_1, \ldots, t_n\) real; in particular, \(q_{jj} = q(a_j)\) for \(1 \leqslant j \leqslant n\) . Moreover, the integral \(\int_{\mathbf{R}} t^n d\gamma(t)\) has the values 1, 0, 1 for \(n = 0, 1, 2\) , respectively (No.~4, Lemma 1). From this, one deduces immediately

\[
\int \dots \int \left(\varepsilon + q \left(\sum_ {j = 1} ^ {n} t _ {j} a _ {j}\right)\right) d \gamma (t _ {1}) \dots d \gamma (t _ {n}) = \varepsilon + \sum_ {j = 1} ^ {n} q (a _ {j}).\tag{59}
\]

Now, the first member of (58) and \(\Phi(0)\) are real numbers; one can therefore replace \(\Phi\) by \(R\Phi\) in the second member of (58). The inequality \(\Phi(0)-\mathcal{R}\Phi\leqslant\varepsilon+q\) and the formulas (56), (58) and (59) then imply

\[
\mu (\mathrm{D}) \leqslant 3 \left(\varepsilon + \sum_ {j = 1} ^ {n} q (a _ {j})\right).\tag{60}
\]

Let us fix the number \(r > 0\) . Since the quadratic form \(h\) is positive, there exist a basis \((e_1, \ldots, e_n)\) of \(T\) and an integer \(m\) between 0 and \(n\) such that

\[
h \left(\sum_ {j = 1} ^ {n} t _ {j} e _ {j}\right) = \sum_ {j = 1} ^ {m} t _ {j} ^ {2}
\]

for \(t_1, \ldots, t_n\) real (Annex, No.~1, Prop. 2). It is then immediate that \(\mathbf{C}_r\) consists of the \(x' \in \mathbf{T}'\) such that

\[
\sum_ {j = 1} ^ {m} \langle e _ {j}, x ^ {\prime} \rangle^ {2} \leqslant r ^ {2}, \quad \sum_ {j = m + 1} ^ {n} \langle e _ {j}, x ^ {\prime} \rangle^ {2} = 0.
\]

For every integer \(l \geqslant 1\) , let \(D_l\) be the set of \(x' \in T'\) satisfying the inequality

\[
\sum_ {j = 1} ^ {m} \left\langle r ^ {- 1} e _ {j}, x ^ {\prime} \right\rangle^ {2} + \sum_ {j = m + 1} ^ {n} \left\langle l e _ {j}, x ^ {\prime} \right\rangle^ {2} > 1.
\]

One sees easily that the sequence \((\mathrm{D}_l)_{l\geqslant 1}\) is increasing with union \(\mathbf{T}' - \mathbf{C}_r\), whence

\[
\mu (\mathrm{T} ^ {\prime} - \mathrm{C} _ {r}) = \lim _ {l \to \infty} \mu (\mathrm{D} _ {l}).\tag{61}
\]

But by (60),

\[
\mu \left(\mathrm{D} _ {l}\right) \leqslant 3 \left(\varepsilon + \sum_ {j = 1} ^ {m} r ^ {- 2} q \left(e _ {j}\right) + \sum_ {j = m + 1} ^ {n} l ^ {2} q \left(e _ {j}\right)\right);\tag{62}
\]

for \(j = m + 1, \ldots, n\) we have \(h(e_j) = 0\) , therefore \(q(e_j) = 0\) . Moreover, \(\operatorname{Tr}(q/h) = \sum_{j=1}^{m} q(e_j)\) (Annex, No.~1, Prop. 2). The relation (55) then follows from (61) and (62).

Q.E.D.

